% GENERATED FILE. Edit canonical lessons, then run npm run generate:books.
% canonical-source-hash: fnv1a64:5c0ae035621525dc
% canonical-lessons: TA-C205-currula, TA-C205-tiraiyarankam, TA-C205-kattitam, TA-C205-mati, TA-C205-kaval-nilaiyam

\chapter{Tour, Cinema, Building, Upper floor, Police station}
\label{ch:ta-tour-cinema-building-upper-floor-police-station}
\hlchaptermodality{205}

\begin{chapteropening}
\textbf{By the end of this chapter:} \emph{I can say a tour, a cinema, a building, an upper floor and a police station.}

Complete the last lesson of chapter 205: I can say a tour, a cinema, a building, an upper floor and a police station.
\end{chapteropening}

\section[\texorpdfstring{\ta{சுற்றுலா} (cuṟṟulā)}{சுற்றுலா (cuṟṟulā)}]{\ta{சுற்றுலா} (cuṟṟulā) — a tour, a sightseeing trip}

\label{lesson:TA-C205-currula}



\begin{quote}
Before the new one: say the Tamil for a supermarket, then the Tamil for a clothes shop.
\end{quote}

\subsection*{You'll want to know: \ta{சுற்றுலா}}
\textbf{\ta{சுற்றுலா}} — \emph{cuṟṟulā} — \textquotedblleft{}a tour, a sightseeing trip\textquotedblright{}.

\begin{cousinweb}
\textbf{\ta{சுற்றுலா}} (\emph{cuṟṟulā}) is native: \emph{cuṟṟu}, \textquotedblleft{}to go around\textquotedblright{} (DEDR 2715), and \emph{ulā}, \textquotedblleft{}a stroll, a procession\textquotedblright{} (DEDR 2693).
\end{cousinweb}

\begin{grammarlens}[title={What you've built}]
One more word for this chapter.
\end{grammarlens}

\subsection*{Your turn}
\begin{itemize}
\raggedright
  \item \emph{Say it:} \emph{cuṟṟulā}
  \item \emph{Say it:} \emph{cuṟṟulā}, once more
  \item \emph{Say it:} say \emph{cuṟṟulā}
  \item \emph{Recall:} say the Tamil for a supermarket, then the Tamil for a clothes shop, then say \emph{cuṟṟulā} again
\end{itemize}

\begin{tcolorbox}[breakable,colback=teal!4,colframe=teal!35!black,title={Before you move on}]
What does \ta{சுற்றுலா} mean? (\textquotedblleft{}A tour, a sightseeing trip\textquotedblright{}.) Say it once more.
\end{tcolorbox}
\section[\texorpdfstring{\ta{திரையரங்கம்} (tiraiyaraṅkam)}{திரையரங்கம் (tiraiyaraṅkam)}]{\ta{திரையரங்கம்} (tiraiyaraṅkam) — a cinema}

\label{lesson:TA-C205-tiraiyarankam}



\begin{quote}
Before the new one: say the Tamil for a clothes shop, then the Tamil for a tour.
\end{quote}

\subsection*{You'll want to know: \ta{திரையரங்கம்}}
\textbf{\ta{திரையரங்கம்}} — \emph{tiraiyaraṅkam} — \textquotedblleft{}a cinema\textquotedblright{}.

\begin{cousinweb}
\textbf{\ta{திரையரங்கம்}} (\emph{tiraiyaraṅkam}) is native \emph{tirai}, \textquotedblleft{}screen\textquotedblright{} (DEDR 3244), and \emph{araṅkam}, \textquotedblleft{}stage, hall\textquotedblright{}, from Sanskrit \emph{raṅga}, \textquotedblleft{}arena, stage\textquotedblright{} (CDIAL 10561), with a vowel put before the \emph{r}.
\end{cousinweb}

\begin{grammarlens}[title={What you've built}]
One more word for this chapter.
\end{grammarlens}

\subsection*{Your turn}
\begin{itemize}
\raggedright
  \item \emph{Say it:} \emph{tiraiyaraṅkam}
  \item \emph{Say it:} \emph{tiraiyaraṅkam}, once more
  \item \emph{Say it:} say \emph{tiraiyaraṅkam}
  \item \emph{Recall:} say the Tamil for a clothes shop, then the Tamil for a tour, then say \emph{tiraiyaraṅkam} again
\end{itemize}

\begin{tcolorbox}[breakable,colback=teal!4,colframe=teal!35!black,title={Before you move on}]
What does \ta{திரையரங்கம்} mean? (\textquotedblleft{}A cinema\textquotedblright{}.) Say it once more.
\end{tcolorbox}
\section[\texorpdfstring{\ta{கட்டிடம்} (kaṭṭiṭam)}{கட்டிடம் (kaṭṭiṭam)}]{\ta{கட்டிடம்} (kaṭṭiṭam) — a building}

\label{lesson:TA-C205-kattitam}



\begin{quote}
Before the new one: say the Tamil for a tour, then the Tamil for a cinema.
\end{quote}

\subsection*{You'll want to know: \ta{கட்டிடம்}}
\textbf{\ta{கட்டிடம்}} — \emph{kaṭṭiṭam} — \textquotedblleft{}a building\textquotedblright{}.

\begin{cousinweb}
\textbf{\ta{கட்டிடம்}} (\emph{kaṭṭiṭam}) is native: \emph{kaṭṭu}, \textquotedblleft{}to build\textquotedblright{} (DEDR 1147), and \emph{iṭam}, \textquotedblleft{}place\textquotedblright{} (DEDR 434): \textquotedblleft{}the built place\textquotedblright{}.
\end{cousinweb}

\begin{grammarlens}[title={What you've built}]
One more word for this chapter.
\end{grammarlens}

\subsection*{Your turn}
\begin{itemize}
\raggedright
  \item \emph{Say it:} \emph{kaṭṭiṭam}
  \item \emph{Say it:} \emph{kaṭṭiṭam}, once more
  \item \emph{Say it:} say \emph{kaṭṭiṭam}
  \item \emph{Recall:} say the Tamil for a tour, then the Tamil for a cinema, then say \emph{kaṭṭiṭam} again
\end{itemize}

\begin{tcolorbox}[breakable,colback=teal!4,colframe=teal!35!black,title={Before you move on}]
What does \ta{கட்டிடம்} mean? (\textquotedblleft{}A building\textquotedblright{}.) Say it once more.
\end{tcolorbox}
\section[\texorpdfstring{\ta{மாடி} (māṭi)}{மாடி (māṭi)}]{\ta{மாடி} (māṭi) — an upper floor, a storey}

\label{lesson:TA-C205-mati}



\begin{quote}
Before the new one: say the Tamil for a cinema, then the Tamil for a building.
\end{quote}

\subsection*{You'll want to know: \ta{மாடி}}
\textbf{\ta{மாடி}} — \emph{māṭi} — \textquotedblleft{}an upper floor, a storey\textquotedblright{}.

\begin{cousinweb}
\textbf{\ta{மாடி}} (\emph{māṭi}) is native Dravidian. DEDR 4796 lists Malayalam \emph{māṭam}, \textquotedblleft{}house with an upper storey\textquotedblright{}. The entry compares Sanskrit and Prakrit words for upper-storeyed houses.
\end{cousinweb}

\begin{grammarlens}[title={What you've built}]
One more word for this chapter.
\end{grammarlens}

\subsection*{Your turn}
\begin{itemize}
\raggedright
  \item \emph{Say it:} \emph{māṭi}
  \item \emph{Say it:} \emph{māṭi}, once more
  \item \emph{Say it:} say \emph{māṭi}
  \item \emph{Recall:} say the Tamil for a cinema, then the Tamil for a building, then say \emph{māṭi} again
\end{itemize}

\begin{tcolorbox}[breakable,colback=teal!4,colframe=teal!35!black,title={Before you move on}]
What does \ta{மாடி} mean? (\textquotedblleft{}An upper floor, a storey\textquotedblright{}.) Say it once more.
\end{tcolorbox}
\section[\texorpdfstring{\ta{காவல்} \ta{நிலையம்} (kāval nilaiyam)}{காவல் நிலையம் (kāval nilaiyam)}]{\ta{காவல்} \ta{நிலையம்} (kāval nilaiyam) — a police station}

\label{lesson:TA-C205-kaval-nilaiyam}



\begin{quote}
Before the new one: say the Tamil for a building, then the Tamil for an upper floor.
\end{quote}

\subsection*{You'll want to know: \ta{காவல்} \ta{நிலையம்}}
\textbf{\ta{காவல்} \ta{நிலையம்}} — \emph{kāval nilaiyam} — \textquotedblleft{}a police station\textquotedblright{}.

\begin{cousinweb}
\textbf{\ta{காவல்} \ta{நிலையம்}} (\emph{kāval nilaiyam}) is native: \emph{kāval}, \textquotedblleft{}guard, watch\textquotedblright{}, from \emph{kā}, \textquotedblleft{}to guard\textquotedblright{} (DEDR 1416), and \emph{nilaiyam}, \textquotedblleft{}station\textquotedblright{}, from \emph{nil}, \textquotedblleft{}to stand\textquotedblright{} (DEDR 3675).
\end{cousinweb}

\begin{grammarlens}[title={What you've built}]
One more word for this chapter.
\end{grammarlens}

\subsection*{Your turn}
\begin{itemize}
\raggedright
  \item \emph{Say it:} \emph{kāval nilaiyam}
  \item \emph{Say it:} \emph{kāval nilaiyam}, once more
  \item \emph{Say it:} say \emph{kāval nilaiyam}
  \item \emph{Recall:} say the Tamil for a building, then the Tamil for an upper floor, then say \emph{kāval nilaiyam} again
\end{itemize}

\begin{tcolorbox}[breakable,colback=teal!4,colframe=teal!35!black,title={Before you move on}]
What does \ta{காவல்} \ta{நிலையம்} mean? (\textquotedblleft{}A police station\textquotedblright{}.) Say it once more.
\end{tcolorbox}
